\documentclass[a4paper]{article}

% Espanol (Mexico) con XeLaTeX: fontspec para fuentes Unicode, polyglossia
% para la division silabica y los nombres del documento en la variante mexicana.
\usepackage{fontspec}
\usepackage{polyglossia}
\setdefaultlanguage[variant=mexican]{spanish}
% polyglossia no traduce los nombres de listings.
\AtBeginDocument{\providecommand{\lstlistingname}{}\renewcommand{\lstlistingname}{Listado}%
  \providecommand{\lstlistlistingname}{}\renewcommand{\lstlistlistingname}{Índice de listados}}
\usepackage{amsmath,amssymb}
\usepackage{graphicx}
\usepackage[margin=2.35cm]{geometry}
\usepackage[most]{tcolorbox}
\usepackage{etoolbox}
\usepackage{listings}
\usepackage{hyperref}
\usepackage{booktabs}
\usepackage{array}
\usepackage{float}
\usepackage{xcolor}

\definecolor{kbBack}{HTML}{F2F6FA}\definecolor{kbFrame}{HTML}{9DB4CC}\definecolor{kbTitle}{HTML}{52708F}
\definecolor{ibBack}{HTML}{FBF5E7}\definecolor{ibFrame}{HTML}{C9A961}\definecolor{ibTitle}{HTML}{7D5F1E}
\definecolor{wbBack}{HTML}{FCF2F0}\definecolor{wbFrame}{HTML}{D6A096}\definecolor{wbTitle}{HTML}{9E4F43}

\tcbset{softbox/.style={enhanced,breakable,boxrule=0.8pt,arc=2pt,left=3mm,right=3mm,top=2mm,bottom=2mm,coltitle=white,fonttitle=\bfseries,toptitle=0.6mm,bottomtitle=0.6mm}}
\newtcolorbox{knowledgebox}[1]{softbox,colback=kbBack,colframe=kbFrame,colbacktitle=kbTitle,title={#1}}
\newtcolorbox{importantbox}[1]{softbox,colback=ibBack,colframe=ibFrame,colbacktitle=ibTitle,title={#1}}
\newtcolorbox{warningbox}[1]{softbox,colback=wbBack,colframe=wbFrame,colbacktitle=wbTitle,title={#1}}

\hypersetup{colorlinks=true,linkcolor=blue!50!black,urlcolor=blue!60!black}
\setlength{\parindent}{2em}
\setlength{\parskip}{0.35em}
\renewcommand{\arraystretch}{1.25}

\newcommand{\notetitle}{Stanford CS329A: 自我改进 AI Agent\\Part 1: 课程总览}
\newcommand{\noteauthors}{AI Course Notes \& Codex}
\newcommand{\notedate}{2026 年 8 月 10 日}
\newcommand{\videochannel}{Stanford Online}
\newcommand{\lecturedate}{2025-09-22}
\newcommand{\videopublishdate}{2026-08-03}
\newcommand{\videoduration}{1:09:42}
\newcommand{\videourl}{https://www.youtube.com/watch?v=6YnLB0XbTnI}

\begin{document}

\begin{titlepage}
\centering
\vspace*{0.8cm}
{\huge\bfseries \notetitle\par}
\vspace{0.6cm}
{\large Aakanksha Chowdhery \& Azalia Mirhoseini\par}
\vspace{0.25cm}
{\large \noteauthors\par}
\vspace{0.8cm}
\includegraphics[width=0.86\textwidth,height=0.44\textheight,keepaspectratio]{cover.jpg}
\vfill
\begin{tcolorbox}[width=0.92\textwidth,colback=black!2!white,colframe=black!55,sharp corners]
\textbf{Curso}: Stanford CS329A Self-Improving AI Agents (Autumn 2025)\par
\textbf{Fecha de la clase}: \lecturedate\quad
\textbf{Fecha de publicación del video}: \videopublishdate\par
\textbf{Fecha de elaboración de las notas}: \notedate\par
\textbf{Canal / duración}: \videochannel\quad / \videoduration\par
\textbf{Enlace del video}: \href{\videourl}{\nolinkurl{\videourl}}\par
\textbf{Normas de elaboración}: \href{https://github.com/wdkns/wdkns-skills}{wdkns/wdkns-skills} (commit \texttt{39f1a04}) y las normas de calidad de este proyecto
\end{tcolorbox}
\end{titlepage}

\tableofcontents
\newpage

\section{Qué intenta responder este curso}

CS329A no se centra en cómo lograr que un modelo de lenguaje responda de forma más humana en una sola llamada, sino en cómo construir un \textbf{sistema capaz de mejorar su desempeño de manera continua mediante cómputo adicional, herramientas externas y retroalimentación del entorno}. El curso enlaza las trayectorias de expansión de capacidad de los últimos años en un solo hilo continuo:

\begin{enumerate}
    \item \textbf{expansión del preentrenamiento (pre-training scaling)}: aumentar parámetros, datos y cómputo de entrenamiento;
    \item \textbf{post-entrenamiento (post-training)}: modificar el comportamiento del modelo mediante datos de instrucción, datos de preferencia y aprendizaje por refuerzo;
    \item \textbf{expansión en tiempo de prueba (test-time scaling)}: con los parámetros fijos, obtener una mayor tasa de éxito mediante muestreo, búsqueda, verificación y razonamiento más extenso;
    \item \textbf{sistemas Agent}: colocar al modelo en un bucle cerrado de acción y retroalimentación, permitiéndole invocar herramientas, leer y escribir en el entorno, planificar y corregirse;
    \item \textbf{automejora}: reconectar generación, verificación, retroalimentación y entrenamiento en un bucle cerrado e iterable.
\end{enumerate}

\begin{importantbox}{La pregunta central de esta sesión}
Cuando el beneficio marginal de simplemente ampliar la escala del preentrenamiento empieza a disminuir, ¿de dónde más puede venir la mejora de capacidad? La respuesta que ofrece esta sesión es: \textbf{trasladar el cómputo de la etapa de entrenamiento a la etapa de inferencia, y elevar la llamada a un solo modelo a un flujo de trabajo de Agent con verificación y retroalimentación del entorno.}
\end{importantbox}

\subsection{Objetivos de aprendizaje}

Al terminar esta sesión, el lector debería poder:

\begin{itemize}
    \item explicar por qué los parámetros, los datos y el cómputo de entrenamiento dan lugar a las scaling laws del preentrenamiento;
    \item distinguir entre zero-shot, few-shot, chain-of-thought e instruction tuning;
    \item explicar cómo RLHF comprime las preferencias humanas en una señal de recompensa optimizable;
    \item deducir la probabilidad de éxito del repeated sampling y comprender la generation--verification gap;
    \item distinguir entre un LLM, un agentic workflow y un Agent autónomo;
    \item comprender todo el curso desde una perspectiva unificada de "generación--verificación--retroalimentación--actualización".
\end{itemize}

\subsection{Resumen de la sección}

El objeto de estudio de CS329A es el sistema completo: el modelo base aporta las capacidades candidatas, el cómputo de inferencia amplía la cobertura de candidatos, el verificador se encarga de seleccionar, las herramientas y el entorno proporcionan retroalimentación real, y el mecanismo de entrenamiento o memoria convierte esa retroalimentación en la capacidad de la siguiente ronda.

\section{Primera vía de escalamiento: Scaling Laws del preentrenamiento}

\subsection{Tres ejes que escalan en conjunto}

La regularidad empírica de la etapa de preentrenamiento es la siguiente: bajo configuraciones razonables de datos y optimización, aumentar la \textbf{cantidad de cómputo de entrenamiento, el volumen de datos de entrenamiento y el número de parámetros del modelo} suele reducir la pérdida de prueba. La forma abstracta más usada es una ley de potencia:

$$
L(x)=L_{\infty}+A x^{-\alpha},\qquad \alpha>0.
$$

\begin{itemize}
    \item $L(x)$: la pérdida de prueba dada una escala de recursos $x$;
    \item $L_{\infty}$: la cota inferior de pérdida difícil de eliminar bajo esa distribución de datos y esos supuestos de modelado;
    \item $A$: una constante de proporcionalidad relacionada con la tarea, los datos y la familia de modelos;
    \item $x$: puede representar cómputo, volumen de datos o número de parámetros;
    \item $\alpha$: el exponente de la ley de potencia, que determina la velocidad a la que baja la pérdida al aumentar los recursos.
\end{itemize}

La importancia de la ley de potencia no radica en que “siempre sea válida”, sino en que en su momento permitió a los equipos predecir, antes de entrenar, el beneficio de un modelo más grande, y con base en eso invertir en cómputo e ingeniería de datos a mayor escala.

\begin{figure}[H]
\centering
\includegraphics[width=0.96\textwidth]{figures/scaling-laws.jpg}
\caption{Scaling laws del preentrenamiento: los tres ejes de escalamiento de cómputo, datos y parámetros.\protect\footnotemark}
\end{figure}
\footnotetext{Intervalo del video: 00:02:50--00:03:53.}

\begin{knowledgebox}{“Más grande” no solo significa más parámetros}
Un modelo con un número de parámetros extremadamente grande, si carece de datos suficientes, terminará memorizando repetidamente un conjunto limitado de muestras; si el cómputo de entrenamiento es insuficiente, los parámetros no se optimizarán del todo; y si la calidad de los datos es mala, el cómputo adicional solo servirá para ajustar el ruido con mayor precisión. Por eso, el núcleo del scaling moderno es la \textbf{configuración coordinada de parámetros, tokens y cómputo efectivo}.
\end{knowledgebox}

\subsection{De BERT y GPT-2 a GPT-4: la escala como fuente de capacidad}

Entre 2,018 y 2,024, la escala de los modelos de lenguaje dominantes creció con rapidez. El curso usa esta historia para subrayar dos puntos: primero, que aumentar la escala trajo consigo una mejora continua del desempeño promedio; segundo, que ciertas capacidades no se vuelven visibles de forma gradual en modelos pequeños, sino que se convierten de golpe en un comportamiento observable al alcanzar cierta escala.

\begin{figure}[H]
\centering
\includegraphics[width=0.92\textwidth]{figures/model-growth.jpg}
\caption{El rápido crecimiento de la escala de parámetros de los modelos de lenguaje grandes.\protect\footnotemark}
\end{figure}
\footnotetext{Intervalo del video: 00:03:55--00:04:46.}

Aquí conviene evitar un malentendido común: que la escala de parámetros crezca en la figura no significa que los parámetros sean la única variable causal. La proporción de los datos, la arquitectura, el optimizador, la estabilidad del entrenamiento y los métodos de post-entrenamiento cambiaron todos en el mismo periodo, de modo que la curva histórica expresa el \textbf{resultado del escalamiento integral de un sistema de ingeniería}.

\subsection{Zero-shot, few-shot y aprendizaje en contexto}

Después de GPT-3, un cambio clave es que el modelo puede entender la tarea directamente a partir del prompt:

\begin{itemize}
    \item \textbf{Zero-shot}: solo se da la descripción de la tarea y la entrada, sin ejemplos;
    \item \textbf{One-shot}: además se da un ejemplo de entrada--salida;
    \item \textbf{Few-shot}: se dan unos cuantos ejemplos para que el modelo infiera del contexto el mapeo de la tarea.
\end{itemize}

Este tipo de capacidad se llama \textbf{in-context learning}: los pesos del modelo no se actualizan, pero el patrón presente en el contexto cambia el comportamiento de ese cómputo hacia adelante en particular.

\begin{figure}[H]
\centering
\includegraphics[width=0.94\textwidth]{figures/zero-few-shot.jpg}
\caption{Diferencia estructural entre los prompts de zero-shot y few-shot.\protect\footnotemark}
\end{figure}
\footnotetext{Intervalo del video: 00:05:40--00:06:38.}

\begin{warningbox}{El aprendizaje en contexto no es aprendizaje de parámetros}
Los ejemplos de few-shot no se escriben de forma permanente en los pesos durante la inferencia. Es más bien como establecer una interpretación temporal de la tarea para la solicitud actual. Una vez que termina el contexto, ese “resultado de aprendizaje” por lo general no se conserva de una solicitud a otra, a menos que el sistema cuente con memoria externa o un mecanismo de entrenamiento continuo.
\end{warningbox}

\subsection{Chain-of-Thought: llevar también el proceso intermedio al contexto}

Lo esencial del prompting de chain-of-thought (CoT) no es escribir unas cuantas frases más, sino aportar una \textbf{trayectoria de razonamiento intermedio que va del problema a la respuesta}. En comparación con mostrar únicamente la respuesta final, el modelo puede imitar la forma de descomponer el problema, actualizar el estado y realizar cálculos locales.

\begin{figure}[H]
\centering
\includegraphics[width=0.96\textwidth]{figures/chain-of-thought.jpg}
\caption{Comparación entre el prompting estándar y el prompting de chain-of-thought.\protect\footnotemark}
\end{figure}
\footnotetext{Intervalo del video: 00:07:18--00:09:01.}

El curso subraya que el efecto del CoT se refuerza claramente al aumentar la escala del modelo. En modelos pequeños, dar ejemplos de razonamiento puede aumentar solo la carga de imitación; en modelos suficientemente grandes, puede desencadenar un cómputo de varios pasos más confiable.

\begin{figure}[H]
\centering
\includegraphics[width=0.92\textwidth]{figures/cot-emergence.jpg}
\caption{La capacidad de chain-of-thought presenta un salto notable en modelos más grandes.\protect\footnotemark}
\end{figure}
\footnotetext{Intervalo del video: 00:09:01--00:10:18.}

\begin{importantbox}{De “predecir el siguiente token” a “ejecutar un proceso de cómputo”}
El CoT muestra que el mismo next-token predictor puede comportarse, dentro del contexto, como un calculador iterativo. El modelo no cambia su función objetivo, pero al generar tokens intermedios, descompone un mapeo difícil en varios mapeos locales más sencillos. Los reasoning model y los Agent posteriores retoman esta misma idea: \textbf{permitir que el sistema consuma más pasos intermedios antes de dar la respuesta final.}
\end{importantbox}

\subsection{Resumen de la sección}

El scaling del preentrenamiento amplía a la vez la capacidad promedio del modelo y el comportamiento potencial que puede activarse mediante el prompt. Zero-shot, few-shot y CoT muestran cómo el contexto puede reorganizar el cómputo del modelo sin actualizar los pesos, pero que el modelo “sepa algo” y que “actúe conforme a la intención humana” siguen siendo dos cosas distintas.
\section{Segunda vía de extensión: Post-Training y las preferencias humanas}

\subsection{Por qué el modelo preentrenado todavía no puede ser un asistente}

El objetivo del preentrenamiento exige que el modelo ajuste la distribución de los textos de internet. Aprende conocimiento, gramática y patrones de razonamiento, pero no tiene una razón natural para seguir instrucciones del usuario, rechazar solicitudes dañinas o preferir respuestas concisas y útiles. La función del post-training es cambiar la \textbf{distribución del comportamiento}: hacer que el modelo, sobre la capacidad de conocimiento que ya tiene, forme una interfaz de tarea más estable.

\begin{figure}[H]
\centering
\includegraphics[width=0.94\textwidth]{figures/alignment.jpg}
\caption{El post-training guía al modelo preentrenado de propósito general hacia la región de comportamiento que se espera de él.\protect\footnotemark}
\end{figure}
\footnotetext{Intervalo de la explicación en video: 00:12:41--00:14:07.}

\subsection{Instruction Tuning: primero aprender a "responder según la instrucción"}

El instruction tuning usa datos de instrucción--respuesta, humanos o sintéticos, para hacer fine-tuning supervisado. Su valor está en establecer una interfaz unificada: ya sea traducción, preguntas y respuestas, resumen o generación de código, todo puede expresarse como "instrucción + contexto $\rightarrow$ respuesta".

Si los datos supervisados son $\mathcal{D}_{\mathrm{SFT}}=\{(x_i,y_i)\}$, el objetivo habitual es:

$$
\mathcal{L}_{\mathrm{SFT}}(\theta)
=-\sum_i \log \pi_{\theta}(y_i\mid x_i).
$$

\begin{itemize}
    \item $x_i$: la instrucción y el contexto de entrada;
    \item $y_i$: la respuesta de demostración;
    \item $\pi_{\theta}$: la política del modelo de lenguaje con parámetros $\theta$;
    \item $\mathcal{L}_{\mathrm{SFT}}$: la pérdida de log-verosimilitud negativa del fine-tuning supervisado.
\end{itemize}

El SFT puede enseñar el formato y el seguimiento de tareas, pero le cuesta mucho representar con una sola "respuesta estándar" las diferencias de calidad entre respuestas abiertas, por lo que también se necesita el aprendizaje de preferencias.

\subsection{RLHF: convertir las preferencias de ranking en una función de recompensa}

El flujo típico de RLHF consta de tres pasos:

\begin{enumerate}
    \item muestrear varias respuestas para el mismo prompt;
    \item pedir a personas que comparen las respuestas y entrenar con ello un modelo de recompensa $r_{\phi}(x,y)$;
    \item usar aprendizaje por refuerzo para aumentar la probabilidad de las respuestas con recompensa alta, a la vez que se limita cuánto se desvía la política del modelo de SFT.
\end{enumerate}

\begin{figure}[H]
\centering
\includegraphics[width=0.95\textwidth]{figures/rlhf.jpg}
\caption{Se entrena un modelo de recompensa con el ranking humano y luego se usa la señal de recompensa para optimizar el modelo de lenguaje.\protect\footnotemark}
\end{figure}
\footnotetext{Intervalo de la explicación en video: 00:16:54--00:19:05.}

El objetivo de optimización puede resumirse de forma abstracta como:

$$
\max_{\theta}\;\mathbb{E}_{y\sim\pi_{\theta}(\cdot\mid x)}[r_{\phi}(x,y)]
-\beta\,D_{\mathrm{KL}}\!\left(\pi_{\theta}\,\|\,\pi_{\mathrm{ref}}\right).
$$

\begin{itemize}
    \item $r_{\phi}(x,y)$: la puntuación de preferencia que el modelo de recompensa asigna a la respuesta;
    \item $\pi_{\theta}$: el modelo de política que se está optimizando;
    \item $\pi_{\mathrm{ref}}$: el modelo de referencia, que suele provenir del SFT;
    \item $D_{\mathrm{KL}}$: el grado en que la política se desvía del modelo de referencia;
    \item $\beta$: el coeficiente que equilibra el aumento de recompensa y la deriva del comportamiento.
\end{itemize}

\begin{warningbox}{El modelo de recompensa es solo un sustituto aproximado de la preferencia}
Si la política encuentra huecos en el modelo de recompensa, puede obtener una puntuación alta pero producir un comportamiento de baja calidad; esto es el reward hacking. La eficacia del RLHF depende de la cobertura de los datos de preferencia, la consistencia de las anotaciones, la capacidad de generalización del modelo de recompensa y las restricciones sobre la política, no del hecho de que "con una recompensa el modelo ya se alinea de manera natural".
\end{warningbox}

\subsection{Resumen de la sección}

El post-training transforma un modelo generativo de propósito general en un asistente interactivo. El SFT establece la interfaz de instrucciones, el modelo de preferencia condensa las comparaciones humanas y el RL optimiza entre la recompensa y la deriva del comportamiento. Es también la primera vez que se introduce de manera explícita una estructura que reaparecerá una y otra vez a lo largo del curso: \textbf{generar candidatos, evaluar candidatos, actualizar el generador según la evaluación}.

\section{Tercera vía de extensión: Inference-Time Scaling}

\subsection{Con los parámetros fijos, aún se puede cambiar más cómputo por tasa de éxito}

El punto de partida de Large Language Monkeys es muy directo: muestrear la misma pregunta de forma independiente muchas veces; mientras el modelo tenga una probabilidad de éxito distinta de cero para esa pregunta, es posible acertar la solución correcta en alguna muestra. El sistema se divide en dos pasos:

\begin{enumerate}
    \item \textbf{Generation}: generar una gran cantidad de candidatos, buscando coverage;
    \item \textbf{Verification}: identificar la solución correcta entre los candidatos, buscando precision.
\end{enumerate}

\begin{figure}[H]
\centering
\includegraphics[width=0.96\textwidth]{figures/language-monkeys.jpg}
\caption{Large Language Monkeys: generación diversa y selección mediante verificación.\protect\footnotemark}
\end{figure}
\footnotetext{Intervalo de la explicación en video: 00:20:03--00:22:43.}

Supongamos que la probabilidad de éxito independiente de una muestra es $p$; la probabilidad de tener al menos un éxito al muestrear $k$ veces es:

$$
P_{\mathrm{solve}}(k)=1-(1-p)^k.
$$

\begin{itemize}
    \item $p$: probabilidad de obtener un candidato correcto en una sola muestra;
    \item $k$: presupuesto de muestreo en el momento de la inferencia;
    \item $(1-p)^k$: probabilidad de fallar las $k$ veces;
    \item $P_{\mathrm{solve}}(k)$: probabilidad de que exista al menos una solución correcta en el conjunto de candidatos.
\end{itemize}

Cuando $p$ es muy pequeño, se necesita un $k$ muy grande para ver una ganancia apreciable; cuando $p=0$, muestrear más no sirve de nada. Esto explica el límite del test-time scaling: puede explotar capacidades que el modelo \textbf{ya posee con baja probabilidad}, pero no puede crear de la nada un algoritmo que el modelo no sabe en absoluto.

\begin{importantbox}{Coverage y precisión final no son lo mismo}
Oracle pass@$k$ supone que podemos saber cuál de los candidatos es correcto, y mide la cobertura de la generación; un sistema real todavía enfrenta el problema de la selección. Si el verifier no logra identificar la respuesta correcta, un conjunto de candidatos por bueno que sea no se traduce en el pass@1 final. Esta diferencia es precisamente el generation--verification gap.
\end{importantbox}

\subsection{Por qué incluso un modelo pequeño puede ser ``amplificado'' con muestreo masivo}

Los experimentos muestran que, para distintos modelos y distintas preguntas, a menudo existe una probabilidad de éxito de cola larga: un modelo de 8B puede fallar la mayoría de las veces en una pregunta, pero acertar en trayectorias muy escasas. Aumentar el número de muestras va recolectando poco a poco esas trayectorias de éxito poco frecuentes, de modo que, bajo una evaluación oracle, el modelo pequeño se acerca o incluso supera el desempeño de un solo intento de un modelo más fuerte.

\begin{figure}[H]
\centering
\includegraphics[width=0.97\textwidth]{figures/repeated-sampling-results.jpg}
\caption{El muestreo repetido amplía la cobertura de candidatos en tareas de matemáticas, código y demostraciones formales.\protect\footnotemark}
\end{figure}
\footnotetext{Intervalo de la explicación en video: 00:23:49--00:27:12.}

Pero la ganancia de cómputo tiene rendimientos decrecientes: una vez cubiertas las preguntas fáciles, el muestreo adicional se consume sobre todo en las preguntas extremadamente difíciles. El diseño del sistema no debe preguntarse solo ``¿muestrear más mejora el resultado?'', sino también ``¿cuánta tasa de éxito marginal se obtiene por cada unidad de token, latencia y energía consumidos?''.

\subsection{El cómputo de entrenamiento y el cómputo de prueba son dos presupuestos distintos}

Los reasoning model incorporan a la inferencia cadenas de pensamiento más largas, búsqueda y autocorrección. El curso usa las curvas de entrenamiento y de tiempo de prueba de o1 para ilustrar que la capacidad puede crecer tanto con el cómputo de entrenamiento como, una vez fijado el modelo, con el test-time compute.

\begin{figure}[H]
\centering
\includegraphics[width=0.9\textwidth]{figures/o1-scaling.jpg}
\caption{Tanto la fase de entrenamiento como la fase de prueba tienen un eje de cómputo escalable.\protect\footnotemark}
\end{figure}
\footnotetext{Intervalo de la explicación en video: 00:29:40--00:32:07.}

\begin{figure}[H]
\centering
\includegraphics[width=0.94\textwidth]{figures/reasoning-models.jpg}
\caption{Los reasoning model consumen más cómputo intermedio mediante descomposición, evaluación, corrección y alternativas.\protect\footnotemark}
\end{figure}
\footnotetext{Intervalo de la explicación en video: 00:32:08--00:33:55.}

\begin{knowledgebox}{Escalamiento paralelo y escalamiento en serie}
El \textbf{escalamiento paralelo} genera de forma independiente varios candidatos; es fácil de paralelizar, pero los candidatos no comparten experiencia entre sí. El \textbf{escalamiento en serie} hace que los pasos posteriores lean los resultados anteriores y realicen revisión, búsqueda o retroceso; puede ser más eficiente, pero si el estado inicial es erróneo, el sesgo también se propaga a los pasos siguientes. Los sistemas reales suelen combinar ambos.
\end{knowledgebox}

\subsection{El presupuesto de inferencia debe adaptarse a la dificultad de la pregunta}

Asignar el mismo presupuesto de tokens a cada pregunta desperdicia cómputo: las preguntas sencillas se resuelven en un solo intento, mientras que las difíciles requieren más exploración. Un controlador más razonable debería decidir si continuar el cómputo con base en la incertidumbre, la puntuación de verificación, la consistencia entre candidatos o señales intermedias de fallo.

El presupuesto de asignación puede escribirse como una optimización con restricciones:

$$
\max_{k_1,\ldots,k_n}\sum_{i=1}^{n} U_i(k_i)
\quad\text{s.t.}\quad
\sum_{i=1}^{n} c_i k_i\le B.
$$

\begin{itemize}
    \item $k_i$: número de pasos de inferencia o de muestras que recibe la pregunta $i$;
    \item $U_i(k_i)$: utilidad esperada que aporta el cómputo adicional;
    \item $c_i$: costo de cada paso;
    \item $B$: presupuesto total de tokens, latencia o dinero.
\end{itemize}

\subsection{Resumen de la sección}

Inference-time scaling traslada el ``más cómputo'' del clúster de entrenamiento a cada solicitud individual. El repeated sampling mejora la coverage, el reasoning/search mejora la calidad de cada trayectoria individual, y el verifier determina si estos candidatos adicionales pueden traducirse en la precisión final. Muchos de los métodos que se verán a lo largo del curso pueden entenderse como mejoras a estas tres partes.
\section{De la llamada a un LLM a los sistemas Agent}

\subsection{El ciclo mínimo del Agent}

Los límites de entrada y salida de una sola llamada a un LLM son fijos; un Agent, en cambio, mantiene un objetivo y un estado dentro de un ciclo: observa el entorno, decide una acción, ejecuta una herramienta, lee la retroalimentación y luego decide el siguiente paso. Su abstracción mínima es:

$$
s_{t+1}=T(s_t,a_t),\qquad a_t\sim\pi_{\theta}(\cdot\mid s_t,g).
$$

\begin{itemize}
    \item $g$: objetivo del usuario;
    \item $s_t$: estado visible en el paso $t$, que incluye el historial, los resultados de las herramientas y el entorno externo;
    \item $a_t$: la acción elegida por el modelo;
    \item $T$: la transición de estado tras la ejecución de la acción por el entorno;
    \item $\pi_{\theta}$: la política del modelo que decide la siguiente acción.
\end{itemize}

\begin{figure}[H]
\centering
\includegraphics[width=0.9\textwidth]{figures/agent-loop.jpg}
\caption{El Agent forma un ciclo cerrado mediante acciones y retroalimentación del entorno.\protect\footnotemark}
\end{figure}
\footnotetext{Intervalo de la explicación en video: 00:42:20--00:43:53.}

\begin{importantbox}{El núcleo de un Agent no es "llamar al modelo unas cuantas veces más"}
Solo cuando las decisiones posteriores leen resultados de ejecución reales y pueden modificar el plan a partir de ellos, el sistema forma un ciclo de retroalimentación cerrado. La generación de texto en varias rondas sin estado del entorno, ejecución de herramientas ni consumo de retroalimentación se parece más a un prompt pipeline que a un Agent completo.
\end{importantbox}

\subsection{Agentic Workflow: escribir estructura confiable en código}

El agentic workflow usa un flujo de control predefinido para organizar los modelos y las herramientas. Sacrifica parte de la autonomía a cambio de previsibilidad, observabilidad y facilidad de depuración. Por ejemplo, en un flujo de generar, luego revisar y reintentar si falla, la lógica de control está definida por el programa, no decidida libremente por el modelo cada vez.

\begin{figure}[H]
\centering
\includegraphics[width=0.94\textwidth]{figures/agentic-workflows.jpg}
\caption{Dos tipos de workflow: generación--evaluación en serie y generación--agregación en paralelo.\protect\footnotemark}
\end{figure}
\footnotetext{Intervalo de la explicación en video: 00:43:53--00:44:31.}

El curso enumera cinco patrones comunes:

\begin{table}[H]
\centering
\small
\begin{tabular}{p{3.2cm}p{5.4cm}p{5.2cm}}
\toprule
\textbf{Patrón} & \textbf{Mecanismo} & \textbf{Escenario de aplicación} \\
\midrule
Prompt chaining & La salida del paso anterior se vuelve la entrada del siguiente & Generación de documentos, extracción de información, transformación por etapas \\
Routing & Primero se determina el tipo de solicitud y luego se asigna a un modelo o herramienta especializados & Distribución de atención al cliente, cascada de modelos, clasificación de riesgo \\
Parallelization & Se generan varios candidatos o se procesan varias subtareas al mismo tiempo & Muestreo diverso, recuperación independiente, análisis por lotes \\
Orchestrator--worker & Un planificador central divide la tarea, varios workers la ejecutan y al final se integra & Modificación de bases de código, investigación profunda, informes complejos \\
Evaluator--optimizer & El generador y el revisor iteran hasta cumplir el criterio o agotar el presupuesto & Revisión de escritura, pruebas de código, optimización de soluciones \\
\bottomrule
\end{tabular}
\caption{Patrones comunes de agentic workflow}
\end{table}

\begin{figure}[H]
\centering
\includegraphics[width=0.94\textwidth]{figures/workflow-patterns.jpg}
\caption{Un workflow se compone de llamadas a LLM, herramientas, búsqueda, verifier, critic y judge.\protect\footnotemark}
\end{figure}
\footnotetext{Intervalo de la explicación en video: 00:44:32--00:49:00.}

\subsection{Compromisos de ingeniería entre workflow y Agent autónomo}

\begin{itemize}
    \item \textbf{Workflow fijo}: alta controlabilidad, fácil de probar, límite de costo claro, pero carece de adaptabilidad ante situaciones nuevas;
    \item \textbf{Agent autónomo}: puede planificar dinámicamente y elegir herramientas, pero la explosión de ramas de comportamiento hace que los errores se acumulen en trayectorias largas;
    \item \textbf{Sistema híbrido}: fija de forma rígida los límites de alto riesgo, los permisos y los criterios de aceptación, y deja la búsqueda abierta al modelo.
\end{itemize}

\begin{warningbox}{A mayor autonomía, menos basta con evaluar solo la respuesta final}
Un sistema de trayectoria larga puede llegar por accidente a un resultado correcto y aun así haber pasado por acciones peligrosas, ciclos ineficaces o costos inaceptables. La evaluación debe registrar a la vez la calidad final, la corrección de las llamadas a herramientas, el número de pasos, la latencia, el costo, los límites de permisos y el comportamiento de recuperación ante fallos.
\end{warningbox}

\subsection{Resumen de la sección}

El Agent extiende el razonamiento hasta convertirlo en un proceso de acción con estado; el agentic workflow escribe en el flujo de control los patrones comunes de descomposición, routing, paralelización y revisión. Ambos requieren los temas posteriores del curso: la planificación decide qué hacer, las herramientas proporcionan retroalimentación del entorno, el verifier determina si es correcto, y el entrenamiento y la memoria hacen que el sistema mejore la próxima vez.
\section{Tres dominios de aplicación representativos}

\subsection{Coding Agent: el repositorio de código como entorno ejecutable}

Las tareas de programación cuentan con retroalimentación natural: el compilador, las pruebas, el análisis estático y los registros de ejecución pueden funcionar como verifier. El valor de un coding agent no es solo completar código, sino mantener un estado de trabajo en constante cambio: buscar en el repositorio, editar múltiples archivos, ejecutar pruebas, leer los fallos, ubicar la causa raíz e iterar.

\begin{figure}[H]
\centering
\includegraphics[width=0.94\textwidth]{figures/coding-agents.jpg}
\caption{Un coding agent opera simultáneamente el shell, el navegador, el editor y el estado del plan.\protect\footnotemark}
\end{figure}
\footnotetext{Intervalo de la explicación en video: 00:52:44--00:54:34.}

\begin{knowledgebox}{Por qué el código es el terreno de prueba ideal para métodos de automejora}
Las tareas de código combinan generación abierta con verificación automática relativamente barata: el espacio de programas candidatos es enorme, pero las pruebas unitarias, la verificación de tipos y los resultados de ejecución permiten filtrar candidatos con rapidez. Esto convierte al código en un punto de encuentro de alto valor para repeated sampling, search, execution feedback y RL.
\end{knowledgebox}

\subsection{Customer Support Agent: la retroalimentación proviene del estado del negocio}

Un agente de atención al cliente necesita acceso a pedidos, tickets, reembolsos y bases de conocimiento, y debe ejecutar acciones reales del negocio. El criterio de éxito normalmente no es que "la respuesta suene razonable", sino si el problema se resolvió, si el estado se actualizó correctamente y si el usuario necesita volver a contactar al servicio.

\begin{figure}[H]
\centering
\includegraphics[width=0.93\textwidth]{figures/customer-support-agents.jpg}
\caption{Un agente de atención al cliente conecta la recuperación de información, el juicio y las acciones del negocio.\protect\footnotemark}
\end{figure}
\footnotetext{Intervalo de la explicación en video: 00:54:34--00:55:43.}

La dificultad de este tipo de sistemas radica en los permisos y las acciones irreversibles: el modelo puede sugerir un reembolso, pero antes de ejecutarlo realmente suele requerirse verificación de reglas, límites de monto o confirmación humana.

\subsection{Research Agent: dividir el proceso de investigación en etapas verificables}

Un agente de investigación abarca la generación de ideas, la búsqueda bibliográfica, el diseño experimental, la ejecución de código, el análisis de resultados y la redacción del artículo. El verifier varía mucho entre etapas: en la búsqueda se puede comprobar si las citas existen, en el código se pueden correr pruebas, pero las conclusiones experimentales requieren rigor estadístico y juicio del dominio.

\begin{figure}[H]
\centering
\includegraphics[width=0.96\textwidth]{figures/research-agents.jpg}
\caption{El flujo de ideas, iteración experimental y redacción del artículo de un research agent.\protect\footnotemark}
\end{figure}
\footnotetext{Intervalo de la explicación en video: 00:56:00--00:59:03.}

\begin{warningbox}{"Completar automáticamente el informe" no equivale a completar la investigación}
Un research agent es especialmente propenso a producir resultados que parecen completos en la superficie pero con evidencia débil en cuanto a la veracidad de las citas, la reproducibilidad experimental y la explicación causal. Un sistema de alta calidad debe incorporar en su flujo de trabajo la verificación de fuentes, los registros de experimentos, las pruebas estadísticas y la búsqueda de contraejemplos, en lugar de optimizar únicamente la fluidez del texto final.
\end{warningbox}

\subsection{Resumen de la sección}

El punto en común de los tres tipos de aplicaciones es que el modelo debe salir del espacio puramente textual e interactuar con un estado externo cambiante. Sus diferencias provienen principalmente del costo y la fiabilidad del verifier: la retroalimentación de la ejecución de código es la más barata, la señal de éxito del negocio llega con más retraso, y la corrección de la investigación es la más difícil de automatizar.
\section{Perspectiva unificada: el ciclo cerrado de automejora}

Esta clase puede resumirse en un ciclo de cuatro etapas:

\begin{enumerate}
    \item \textbf{Generate}: el modelo o varios Agents producen acciones, soluciones o planes candidatos;
    \item \textbf{Verify}: pruebas, un reward model, reglas o señales del entorno evalúan a los candidatos;
    \item \textbf{Select / Act}: se selecciona un candidato y se ejecuta en el entorno;
    \item \textbf{Update}: la retroalimentación se usa para corregir la trayectoria actual, actualizar la memoria o entrenar los parámetros.
\end{enumerate}

\begin{importantbox}{Mapa conceptual de todo el curso}
El test-time scaling mejora principalmente Generate y Select; la robust verification mejora Verify; el tool use y el planning mejoran Act; el RL y los synthetic data reescriben la retroalimentación en los parámetros; la memory y el continual learning conservan la experiencia entre tareas. Lo que llamamos self-improving agent es precisamente lograr que esta cadena pueda repetirse de manera estable y medible.
\end{importantbox}

\subsection{Lista de verificación al diseñar un sistema de automejora}

\begin{itemize}
    \item ¿De dónde viene la diversidad de los candidatos? ¿No se está simplemente repitiendo el mismo error?
    \item ¿El verifier es más confiable que el generator, o comparten el mismo punto ciego?
    \item ¿La retroalimentación es inmediata, está diferida, o solo puede aproximarse con una métrica proxy?
    \item ¿Qué acciones son reversibles y cuáles requieren permisos o confirmación humana?
    \item ¿El beneficio marginal del cómputo adicional compensa el costo y la latencia?
    \item ¿La actualización podría provocar reward hacking, olvido o deriva de las capacidades?
\end{itemize}

\subsection{Resumen de la sección}

Al separar el modelo, el verificador, las herramientas y el mecanismo de actualización, es posible ubicar con mayor claridad los cuellos de botella del sistema. Muchos fallos que parecen deberse a que "el modelo no es suficientemente inteligente" en realidad provienen del verificador, la gestión de estado, la interfaz de herramientas o la asignación del presupuesto, y no necesariamente requieren volver a entrenar un modelo base más grande.

\section{Lecturas adicionales}

\begin{itemize}
    \item \href{https://arxiv.org/abs/2407.21787}{Brown et al., \emph{Large Language Monkeys: Scaling Inference Compute with Repeated Sampling}}
    \item \href{https://www.anthropic.com/research/building-effective-agents}{Anthropic, \emph{Building Effective Agents}}
    \item \href{https://cs329a.stanford.edu/}{Página del curso Stanford CS329A y reading list}
    \item \href{https://www.youtube.com/playlist?list=PLangBM27OtEA}{Lista de reproducción oficial de 9 partes}
\end{itemize}

\section{Síntesis y ampliación}

Esta clase partió de las scaling laws del preentrenamiento para explicar por qué las capacidades de los modelos grandes dependieron durante mucho tiempo de la expansión conjunta de parámetros, datos y cómputo de entrenamiento; después mostró, mediante zero/few-shot y chain-of-thought, cómo el contexto activa el cómputo latente interno del modelo; explicó, mediante instruction tuning y RLHF, cómo el post-entrenamiento moldea el comportamiento; luego, con Large Language Monkeys y los reasoning models, introdujo el cómputo en tiempo de inferencia; y finalmente situó estas capacidades dentro de un sistema de Agent con herramientas y retroalimentación.

Los tres juicios que más vale la pena conservar son:

\begin{enumerate}
    \item \textbf{La capacidad no reside solo en los pesos, sino también en el proceso de cómputo que se puede invocar.} El mismo modelo puede mostrar tasas de éxito muy distintas según el muestreo, la búsqueda y el presupuesto de razonamiento.
    \item \textbf{La capacidad de generación y la capacidad de verificación deben medirse por separado.} Que exista una "respuesta correcta" entre los candidatos no equivale a que el sistema "pueda entregar la respuesta correcta".
    \item \textbf{El valor de un Agent proviene del ciclo cerrado.} Solo con ejecución real, consumo de retroalimentación y actualización de estado se puede elevar un modelo de lenguaje a un sistema capaz de completar tareas de manera sostenida.
\end{enumerate}

La siguiente clase se dedicará específicamente al test-time compute scaling: por qué el muestreo repetido presenta rendimientos con ley de potencias, cómo se comparan el muestreo paralelo y la revisión secuencial, cómo el reward model guía la búsqueda, y en qué casos el preentrenamiento adicional sigue siendo mejor que el cómputo adicional en la inferencia.

\end{document}
